\section{The layer sweep}
\label{app:layers}

Activations were re-extracted at six layers for the 80-trial subset that permits it,
which carries 7 events. Raw AUC for $V_{\text{int}}$[desperate] against reward
hacking varies widely across layers, and residualizing on response length removes
the variation along with the advantage.

{\centering\footnotesize
\begin{tabular}{@{}lcccc@{}}
\toprule
Layer & raw AUC & length-resid.\ AUC & 95\% CI & $\rho$(probe, length) \\
\midrule
13 & 0.691 & 0.333 & $[0.194, 0.486]$ & 0.619 \\
26 & 0.450 & 0.262 & $[0.118, 0.419]$ & 0.445 \\
\textbf{39} & \textbf{0.918} & \textbf{0.646} & $\mathbf{[0.384, 0.877]}$ & 0.603 \\
52 & 0.708 & 0.532 & $[0.282, 0.781]$ & 0.392 \\
53 (steered) & 0.677 & 0.524 & $[0.279, 0.773]$ & 0.373 \\
65 & 0.472 & 0.360 & $[0.113, 0.622]$ & 0.344 \\
\bottomrule
\end{tabular}\par}

\vspace{4pt}
Layer 39's apparent superiority is the same length artifact seen at layer 53, and
its residualized interval spans chance. Length in tokens alone reaches AUC $0.912$
on this subset, and the probe correlates with length at every layer.

Layers 13 and 26 have residualized intervals falling below chance. We read the sign
rather than the significance, since these are 7 events and the sweep makes 6 by 50
comparisons without correction. The sign reflects different geometry rather than a
flip of the same axis. Comparing the full 50-direction residualized-AUC profile
across layers, 13 and 26 are nearly uncorrelated with layer 53 ($\rho = 0.24$ for
both) and with each other ($\rho = 0.04$), while layers 39, 52, 53 and 65 form a
tightly correlated block (Spearman $0.61$ to $0.99$). The behaviorally associated
axis is a late-layer phenomenon. Layer 39 correlates with the steered layer 53 at
$\rho = 0.82$, so the two sites largely measure the same thing, which is a further
reason that relocating the intervention there would not have tested anything new.
